\documentclass[12pt,a4paper]{article}
\usepackage[utf8]{inputenc}
\usepackage[brazil]{babel}
\usepackage{graphicx}
\usepackage{amssymb, amsfonts, amsmath}
\usepackage{color}
\usepackage{float}
\usepackage{enumerate}
%\usepackage{subfigure}
\usepackage[top=2.5cm, bottom=2.5cm, left=1.25cm, right=1.25cm]{geometry}

\newcommand{\sen}{\,\textrm{sen}\,}
\newcommand{\tg}{\,\textrm{tg}\,}

\begin{document}
\pagestyle{empty}

\begin{center}
\begin{tabular}{ccc}
\begin{tabular}{c}
\includegraphics[scale=0.25]{../../biblioteca/imagem/brasao-de-armas-brasil} \\
\end{tabular} & 
\begin{tabular}{c}
Ministério da Educação \\
Universidade Federal dos Vales do Jequitinhonha e Mucuri \\
Faculdade de Ciências Sociais, Aplicadas e Exatas - FACSAE \\
Departamento de Ciências Exatas - DCEX \\
Disciplina: Cálculo Numérico\\
Prof.: Luiz C. M. de Aquino\\
\end{tabular} &
\begin{tabular}{c}
\includegraphics[scale=0.25]{../../biblioteca/imagem/logo-ufvjm} \\
\end{tabular}
\end{tabular}

Aluno(a): \rule{0.5\textwidth}{0.01cm} \quad Data: \rule{0.6cm}{0.01cm} / \rule{0.6cm}{0.01cm} / \rule{1.25cm}{0.01cm}
\end{center}

\begin{center}
 \textbf{Avaliação I -- 2$^a$ Chamada}
\end{center}

\textbf{Instruções}
\begin{itemize}
 \item Todas as justificativas necessárias na solução de cada questão devem estar presentes nesta avaliação;
 \item As respostas finais de cada questão devem estar escritas de caneta;
 \item Esta avaliação tem um total de 25,0 pontos.
\end{itemize}

\begin{enumerate}
 \item \textbf{[6,0 pontos]} Seja $x$ um número natural qualquer. Considere que $n$ seja um quadrado perfeito mais próximo de $x$. Prove que $\sqrt{x}\approx \dfrac{x+n}{2\sqrt{n}}$. 
 (Observação: dizemos que $n$ é um quadrado perfeito se existe um natural $m$ tal que $n = m^2$.)
 \item \textbf{[4,5 pontos]} Use o Método da Secante para encontrar a raiz aproximada da função definida por $f(x) = \cos x - \dfrac{1}{5}$ no intervalo $[1;\,2]$ (considere uma tolerância de $10^{-5}$).
 \item \textbf{[4,5 pontos]} Utilize o Método de Newnton para determinar uma aproximação para a raiz da função polinomial definida por $p(x) = 2x^4 -2x^3 -22x^2 - 10x + 8$ no intervalo $[0;\,1]$ (considere uma tolerância 
 de $10^{-5}$).
 \item \textbf{[5,0 pontos]} A cada passo no Método da Falsa Posição, escolhemos $x_k = \dfrac{a_k|f(b_k)| + b_k|f(a_k)|}{|f(a_k)|+|f(b_k)|}$, sendo que no intervalo $[a_k;\,b_k]$ temos $f(a_k)f(b_k)<0$. Prove que 
 esta escolha de $x_k$ coincide com a abscissa do ponto de interseção entre o eixo $x$ e a reta passando por $(a_k,\,f(a_k))$ e $(b_k,\,f(b_k))$.

 \item \textbf{[5,0 pontos]} Seja $f:[a;\, b]\to [a;\,b]$ uma função contínua em todo o seu domínio. Prove que o gráfico de $f$ e de $g(x) = x$ tem pelo 
menos um ponto em comum.

\end{enumerate}
\end{document}
